\chapter{Experimental Methodology}
\label{ch:methodology}

\section{Dataset Description}
\label{sec:dataset-description}

All experiments reported in this thesis are conducted on the CICIDS2018 dataset, specifically the NetFlow v3 variant published by the Canadian Institute for Cybersecurity \autocite{Sharafaldin2018}. The dataset was generated in a controlled laboratory environment over ten days in February and March 2018, during which a diverse set of attack scenarios was executed against a realistic enterprise network topology comprising victim hosts, a gateway router, and a dedicated attacker infrastructure. Traffic was captured as full packet captures and subsequently converted to bidirectional network flow records using the CICFlowMeter tool, yielding structured telemetry records that capture connection-level statistics without retaining raw packet payloads.

The NetFlow v3 representation used in this work contains 53 features per flow record, encompassing protocol-level attributes (protocol number, source and destination ports, TCP flag counts), volumetric measures (total forward and backward packets, total bytes transmitted in each direction, packet length statistics), timing characteristics (flow duration, inter-arrival times, active and idle durations), and derived rate features (packets per second, bytes per second). This feature set is substantially richer than the PCAP-derived features available in earlier CICIDS editions and provides the LLM agents with sufficient numerical detail to reason meaningfully about protocol conformance, statistical anomalies, and temporal behaviour. Source and destination IP addresses are preserved as private-range addresses in the 172.31.x.x subnet, reflecting the anonymisation applied to the laboratory infrastructure. This anonymisation is consequential for the MCP comparison experiments described in Section~\ref{sec:mcp-comparison-protocol}, as it renders external threat intelligence lookups against public IP reputation services uninformative by design.

The complete dataset comprises 20.12 million flow records spanning fourteen labelled attack categories alongside a dominant benign class. The fourteen attack types are: FTP-BruteForce, SSH-Bruteforce, DoS-SlowHTTPTest, DoS-Slowloris, DoS-Hulk, DoS-GoldenEye, DDoS-LOIC-HTTP, DDoS-LOIC-UDP, DDoS-HOIC, Bot, Infiltration, Brute\_Force-Web, Brute\_Force-XSS, and SQL\_Injection. The class distribution is heavily skewed, with approximately 87 per cent of flows labelled benign and the remaining 13 per cent distributed unevenly across the fourteen attack categories. This skew is not an artefact of dataset construction but rather reflects the empirical reality of enterprise network traffic, where the vast majority of connections are legitimate. It is precisely this skew that makes na\"{i}ve evaluation on balanced datasets misleading: a system that performs well at 50/50 class balance may exhibit catastrophic false positive rates when confronted with the true traffic distribution, as observed in this work's own pilot experiments with GPT-4o-mini (discussed in Chapter~5).

For the purposes of this research, the dataset is partitioned into three non-overlapping chronological splits. The development partition contains 7.04 million flows and is used for Random Forest training, feature selection, and hyperparameter calibration. The validation partition contains 5.03 million flows and is used for threshold selection, pilot experiments, and intermediate evaluation during system development. The test partition contains 8.05 million flows and is held out for final evaluation, remaining entirely unseen until the concluding experiments reported in Chapter~6. This three-way split follows established practice in machine learning research \autocite{Cawley2010} and is essential for obtaining unbiased performance estimates on data that genuinely represents novel traffic.

A subset of 14 features is used as the input feature set for the Random Forest tier-one pre-filter, selected from the full 53-feature set through a combination of variance thresholding and pairwise correlation analysis. One candidate feature, \texttt{FLOW\_DIRECTION}, was found to be absent from the CICIDS2018 NetFlow v3 schema despite its presence in earlier versions of the dataset; accordingly, the final feature set comprises 14 rather than the originally targeted 15 features. These 14 features include protocol number, source and destination port, packet counts in each direction, byte totals, flow duration, TCP flag fields, and several packet-length statistics. The full feature list is recorded in \texttt{data/feature\_sets/stage1\_features.json} to ensure reproducibility.

\section{Batch Construction}
\label{sec:batch-construction}

Each Stage 1 experiment evaluates AMATAS on a batch of 1,000 flow records drawn from the validation partition. The batch composition is fixed at 50 attack flows and 950 benign flows, yielding a five per cent attack prevalence. This composition is deliberately chosen to approximate realistic production conditions. Empirical surveys of enterprise network traffic find attack prevalence in the range of one to thirteen per cent during active incident periods \autocite{Sommer2010}, and five per cent represents a conservative but plausible operating point that is considerably more demanding than the balanced benchmarks commonly used in intrusion detection literature. Using a balanced dataset would inflate recall and precision estimates and obscure the false positive behaviour that is the primary practical concern for operational deployments.

Attack flows for each experiment are sampled from the validation partition, restricted to the specific attack type under evaluation. The 50 attack flows are sampled without replacement using a fixed random seed of 42, applied consistently across all experiments to ensure that any given attack type's evaluation is reproducible by any researcher with access to the same dataset partition. The 950 benign flows are drawn from the same partition using the same seed, stratified to avoid temporal concentration in any single time window. The precise flows selected, along with their ground truth labels, are serialised to a \texttt{ground\_truth.json} file within each batch directory, allowing retrospective verification of any reported result.

Flow ordering within each batch follows the chronological-by-source-IP convention. Flows are first sorted by source IP address and then, within each source IP group, by \texttt{FLOW\_START\_MILLISECONDS} in ascending order. This ordering serves two purposes. First, it groups flows from the same source together, enabling the Temporal Agent to observe the sequence of connections originating from any given address without requiring a separate indexing step. Second, it preserves the temporal structure of the original traffic, which is essential for detecting attack patterns that manifest across multiple flows over time, such as the progressive connection escalation characteristic of brute-force attacks and the sustained high-rate behaviour of Denial of Service floods. Randomising flow order within a batch would artificially disrupt this temporal signal and would misrepresent the information available to a real-time detection system.

Each batch directory contains three files: \texttt{flows.json}, which stores the 14-feature records for all 1,000 flows in the ordered sequence; \texttt{ground\_truth.json}, which stores the binary label and attack type for each flow; and \texttt{batch\_metadata.json}, which records the batch creation parameters including attack type, sample counts, random seed, partition source, and creation timestamp. The metadata file enables automated auditing of batch provenance and ensures that any result file can be traced back to its exact input data.

\section{Evaluation Metrics}
\label{sec:evaluation-metrics}

The primary evaluation metric throughout this work is detection recall, also referred to as sensitivity or the true positive rate. Recall is defined as the proportion of attack flows that the system correctly classifies as malicious or suspicious:

\begin{equation}
\text{Recall} = \frac{\text{TP}}{\text{TP} + \text{FN}}
\label{eq:recall}
\end{equation}

\noindent where TP denotes true positives (attack flows correctly flagged) and FN denotes false negatives (attack flows that pass through the system undetected). Recall is prioritised over precision or accuracy for a security-critical reason: an undetected attack is categorically more harmful than an investigated false alarm. Security operations teams can absorb a manageable false positive rate; they cannot recover from threats that evade detection entirely. This prioritisation is consistent with the design philosophy of the tier-one Random Forest pre-filter, which is calibrated at a threshold where recall on the development partition is one hundred per cent, accepting higher false positive rates at that tier as the acceptable cost of eliminating false negatives.

The false positive rate (FPR) is the secondary metric and is defined as the proportion of benign flows incorrectly classified as malicious:

\begin{equation}
\text{FPR} = \frac{\text{FP}}{\text{FP} + \text{TN}}
\label{eq:fpr}
\end{equation}

\noindent where FP denotes false positives (benign flows flagged as malicious) and TN denotes true negatives (benign flows correctly passed). In a production environment processing hundreds of thousands of flows per hour, even a one per cent false positive rate would generate thousands of spurious alerts per day, rendering the system operationally impractical. The Stage 1 experiments therefore aim for an FPR below five per cent as the acceptable operating bound, with lower values preferred.

The F1 score is reported as a composite measure of precision and recall, providing a single-figure summary that captures the balance between the two:

\begin{equation}
\text{F1} = \frac{2 \times \text{Precision} \times \text{Recall}}{\text{Precision} + \text{Recall}}
\label{eq:f1}
\end{equation}

\noindent Given the class imbalance in each batch (950 benign versus 50 attack), accuracy is reported as a secondary metric only and is not used as the basis for any substantive comparison, since a system that classifies every flow as benign would achieve 95 per cent accuracy whilst being entirely useless for detection purposes.

Two cost metrics are additionally reported for every experiment. The cost per flow is the total USD expenditure on OpenAI API calls divided by the number of flows processed by the LLM pipeline. Since the tier-one pre-filter eliminates a substantial fraction of flows before they reach the LLM tier, the cost per flow for the full system is considerably lower than the cost per LLM-analysed flow; both quantities are reported separately for transparency. The cost per true positive is the total expenditure divided by the number of attacks correctly detected, providing an economically interpretable measure of detection efficiency that enables direct comparison between system configurations with different costs and different recall rates.

A hard cost limit of five US dollars per batch is imposed on all Stage 1 experiments. This limit is enforced programmatically by the pipeline controller, which checks cumulative expenditure after each flow and terminates the batch with a partial result if the limit is reached. The five-dollar threshold was selected based on pilot experiment cost observations: at the tier-one pre-filter's typical pass-through rate of approximately five per cent, a 1,000-flow batch will present roughly 50 flows to the LLM tier. At the observed cost of approximately 2.8 US cents per LLM-analysed flow, the expected cost per batch is approximately 1.40 dollars, well within the five-dollar limit and providing headroom for batches with higher pass-through rates.

\section{Stage 1 Experimental Protocol}
\label{sec:stage1-protocol}

Stage 1 is the systematic per-attack-type evaluation of the AMATAS v2 two-tier architecture. It comprises fourteen independent experiments, one per attack category in the CICIDS2018 dataset, each conducted on a 1,000-flow batch constructed according to the protocol described in Section~\ref{sec:batch-construction}. The purpose of Stage 1 is to characterise detection performance as a function of attack type rather than as an aggregate figure, since different attack categories present fundamentally different challenges at the flow level: volumetric denial-of-service attacks produce extreme packet-rate and byte-volume anomalies that are numerically conspicuous, whilst infiltration and web application attacks may produce individual flows that are statistically unremarkable and require contextual or behavioural reasoning to identify.

Each Stage 1 experiment proceeds through the following sequence. First, the 1,000 flows in the batch are passed through the tier-one Random Forest pre-filter. The pre-filter computes a maliciousness probability for each flow using the trained Random Forest model and applies a threshold of 0.15: flows with predicted maliciousness probability below 0.15 receive an automatic BENIGN verdict without LLM involvement. Flows at or above the threshold are forwarded to the six-agent LLM pipeline. The 0.15 threshold was selected because it achieves one hundred per cent recall on the development partition --- that is, no attack flow in the 7.04-million-flow development set receives a probability below 0.15 --- whilst still diverting approximately 95 per cent of benign flows away from the LLM tier.

Flows that pass the tier-one filter enter the six-agent pipeline, which processes them sequentially in the ordered batch sequence. For each flow, the Protocol Agent, Statistical Agent, Behavioural Agent, and Temporal Agent are invoked in parallel, each receiving the full 14-feature flow record formatted as a structured natural language prompt. The Temporal Agent additionally receives the preceding flows from the same source IP address within the batch, providing inter-flow context that enables the identification of scanning patterns, connection escalations, and rate anomalies visible only across multiple flows. Upon completion of the four parallel calls, their combined outputs are passed to the Devil's Advocate Agent, which receives all four specialist analyses and is prompted to construct the strongest available argument that the flow is benign, regardless of the specialists' verdicts. Finally, the Orchestrator Agent receives all five preceding analyses and synthesises a final verdict, confidence score, attack type prediction, and full reasoning narrative.

All six agents use GPT-4o (OpenAI) as the underlying language model. This model selection reflects the findings of the pilot validation experiments conducted on the CICIDS2018 validation partition: GPT-4o-mini was rejected due to systematic over-prediction of maliciousness that resulted in a false positive rate approaching one hundred per cent on the pilot batch, whilst Claude Sonnet 3.5 was rejected for Stage 1 due to rate limits that proved incompatible with the throughput requirements of 1,000-flow batches under the project timeline. GPT-4o exhibited the most balanced behaviour in pilot experiments, achieving 50--84 per cent recall with false positive rates of zero to five per cent across attack types.

The Devil's Advocate Agent is assigned a weight of 30 per cent in the Orchestrator's consensus calculation. This weighting was calibrated during Phase 3 validation experiments: a weight of 50 per cent was found to be excessively aggressive, suppressing detection of genuine attacks by over-correcting towards benign interpretations. At 30 per cent, the Devil's Advocate applies meaningful pressure against over-eager malicious classifications whilst preserving the sensitivity of the specialist consensus.

Experiment execution is governed by a pipeline control file at \texttt{results/stage1/control.json}. The pipeline polls this file before processing each flow and responds to three possible commands: \texttt{run} instructs the pipeline to proceed normally; \texttt{pause} causes the pipeline to suspend execution and poll again every thirty seconds; and \texttt{stop} causes the pipeline to abort, serialise the partial results with a status field of \texttt{interrupted}, and exit. This control mechanism enables the operator to halt or suspend long-running experiments without losing accumulated results, which is practically important given that a 1,000-flow batch may require several hours to complete when LLM API response times are variable. Live progress is additionally written to \texttt{results/stage1/live\_status.json} every ten flows, enabling real-time monitoring of throughput, cost accumulation, and running accuracy.

Results from each experiment are serialised to a dedicated JSON result file in \texttt{results/stage1/}, named by attack type. Each result file contains the full per-flow audit trail: the tier-one decision and probability, the verdict and reasoning text from each of the six agents, the token counts and costs for each agent call, and the final Orchestrator verdict against the ground truth label. This granular output is the primary artefact of the Stage 1 evaluation and constitutes the explainability data that the thesis argues is AMATAS's principal contribution over traditional classifiers. Aggregate metrics are additionally appended to \texttt{results/stage1/running\_summary.json} upon completion of each experiment.

\section{MCP Comparison Protocol}
\label{sec:mcp-comparison-protocol}

Separate from the AMATAS Stage 1 evaluation, a set of three single-agent experiments is conducted to contextualise the multi-agent architecture against simpler LLM-based approaches. These experiments employ a Model Context Protocol (MCP) architecture in which a single LLM agent receives a network flow and produces a verdict, with varying levels of prompt engineering and tool access. The comparison is motivated by the research narrative's progression from Phase 1 (single-agent MCP) through Phase 2 (prompt-engineered single agent) to Phase 3 (multi-agent AMATAS): quantifying the performance of these simpler configurations isolates the contribution of each architectural increment.

Three configurations are evaluated. Configuration A (zero-shot) provides the agent with a minimal system prompt and the raw flow features, making no attempt to inject domain knowledge or attack signatures. The model used for this configuration is GPT-4o-mini, selected to test whether a cheaper, smaller model can compensate for the absence of prompt engineering through sheer statistical pattern recognition. Configuration B (engineered prompt) provides a detailed system prompt containing definitions of each feature, descriptions of attack-specific signatures, guidance on interpreting TCP flag combinations, and instructions for structured output formatting. The model for this configuration is GPT-4o. Configuration C (engineered prompt plus MITRE tool) adds an MCP-accessible tool that the agent may invoke to retrieve MITRE ATT\&CK technique descriptions by identifier; the model remains GPT-4o. The MITRE tool is included to test whether external structured knowledge improves single-agent reasoning, and to document its limitations when applied to anonymised laboratory traffic.

The MCP comparison batch comprises 100 flows drawn from the validation partition: 30 attack flows and 70 benign flows, yielding a 30 per cent attack prevalence. This composition differs from the Stage 1 batch design for two reasons. First, the MCP comparison is primarily qualitative in intent --- it aims to demonstrate that single-agent approaches face structural limitations that the multi-agent architecture is designed to overcome --- and the higher attack prevalence provides sufficient attack samples for meaningful per-configuration recall comparison at a fraction of the cost of a full 1,000-flow batch. Second, the cost differential between configurations (Configuration A costs approximately 0.01 US dollars versus 0.45 US dollars for Configuration C) makes a matched 1,000-flow comparison economically unjustifiable for what is intended as a contextual rather than primary result.

Each flow in the MCP batch is processed by all three configurations independently, ensuring that any performance differences are attributable to the configuration rather than to sampling variation. The same ground truth labels are used for all three evaluations. Results are reported as recall, FPR, F1, and total cost per 100 flows.

The expected finding motivating this comparison design is that external tool access does not meaningfully improve single-agent performance on this dataset. Because all IP addresses in CICIDS2018 NetFlow v3 are anonymised to private-range addresses, no external threat intelligence service --- whether AbuseIPDB, AlienVault OTX, or any geolocation provider --- can return useful information about the traffic sources. A MITRE ATT\&CK tool that returns technique descriptions conditioned on attacker-supplied identifiers is additionally subject to the limitation that the single agent must first correctly hypothesise the relevant technique before invoking the tool; if the agent's initial hypothesis is incorrect, the tool retrieval will reinforce the error rather than correct it. This expected finding is not a trivial or negative result: it motivates the multi-agent approach by demonstrating that enriching a single-agent system with external knowledge is insufficient and that the architectural diversity of the AMATAS pipeline is a necessary rather than optional design choice.

\section{Data Integrity and Leakage Mitigation}
\label{sec:data-integrity}

A data integrity concern was identified during the preparation of Stage 1 experiments that required methodological resolution before results could be considered valid for thesis reporting. Seven of the fourteen attack types --- specifically those attack categories whose flows were concentrated in the portion of the original CICIDS2018 dataset that was allocated to the development partition during the initial three-way split --- are present in the same data from which the Random Forest tier-one model was trained. This within-split overlap creates a potential for data leakage: if the Random Forest was trained on flows from the same attack-type distribution that it is subsequently evaluated upon within the Stage 1 batches, its performance figures would be optimistic estimates that do not reflect generalisation to genuinely novel traffic.

The concern is specifically relevant to the tier-one pre-filter's recall measurement. If the Random Forest has seen flows from the same source distribution during training, its ability to correctly assign elevated maliciousness probabilities to attack flows drawn from the same distribution may overstate what the model would achieve on traffic from a truly held-out attack campaign. For the LLM pipeline components, which do not train on historical flow records and reason purely from the feature values presented at inference time, the leakage concern does not apply in the same way: a language model's behaviour is not altered by the composition of the evaluation batch, since it has no learning mechanism that adapts to observed data.

The resolution adopted in this work is a re-partitioning of the development split. The 7.04-million-flow development partition is further divided into a training subset (80 per cent, approximately 5.63 million flows) and a held-out evaluation subset (20 per cent, approximately 1.41 million flows), using stratified random sampling with seed 42 to preserve class distribution across the split. The Random Forest is retrained exclusively on the 80 per cent training subset. The 0.15 probability threshold is then re-validated on the held-out 20 per cent subset to confirm that one hundred per cent attack recall is maintained at this threshold for all fourteen attack types. The re-validation confirms this property: at threshold 0.15, the retrained Random Forest flags all attack flows in the held-out subset as at or above the threshold, whilst the false positive rate on benign flows in the held-out subset remains below ten per cent, consistent with the original model's behaviour.

Stage 1 evaluation batches for the seven potentially affected attack types are drawn from the validation partition, which is entirely disjoint from both the training subset and the held-out evaluation subset of the development partition. For the remaining seven attack types, which were not present in the development partition in meaningful quantities, the validation partition is likewise used as the source. This ensures that, for all fourteen attack types, the flows presented to the AMATAS system during Stage 1 evaluation were not seen --- either by the Random Forest during training or as part of threshold calibration --- prior to the evaluation run. The result is a methodologically clean experimental design in which all reported performance figures reflect genuine generalisation rather than in-sample memorisation.

The batch provenance metadata stored in \texttt{batch\_metadata.json} records the source partition for every flow in every batch, enabling independent verification that no evaluation flow originates from the development partition training subset.

\section{Reproducibility Measures}
\label{sec:reproducibility}

Reproducibility is treated as a first-class methodological requirement throughout this work. Network intrusion detection research has a well-documented reproducibility problem: many published results cannot be replicated because the evaluation data, preprocessing code, or model weights are not released alongside the paper \autocite{Kenyon2020, Ring2019}. AMATAS addresses this through several interlocking mechanisms.

All random sampling operations use a global seed of 42, applied consistently through Python's \texttt{random} module and NumPy's random number generator. This seed governs batch construction, the development partition's 80/20 internal split, the Random Forest's bootstrap sampling during training, and any feature selection steps that involve random tie-breaking. Given the same seed and the same source CSV files, any researcher can reconstruct the exact set of flows in each evaluation batch and verify that the ground truth labels match those stored in the published \texttt{ground\_truth.json} files.

Model versioning is maintained at two levels. The Random Forest model is serialised to \texttt{models/tier1\_rf.pkl} using Python's \texttt{pickle} module following scikit-learn 1.3 conventions, with the scikit-learn version recorded in the batch metadata. The LLM agents call the OpenAI API using model identifier \texttt{gpt-4o}, pinned to the \texttt{gpt-4o-2024-08-06} snapshot for all Stage 1 experiments. Pinning to a specific snapshot version rather than the floating \texttt{gpt-4o} alias is essential because OpenAI periodically updates the models behind floating aliases, which can alter response behaviour and invalidate previously obtained results. The model snapshot identifier is recorded in every result file's metadata block.

Feature sets are stored as versioned JSON files in \texttt{data/feature\_sets/}. The 14-feature set used throughout Stage 1 is stored in \texttt{stage1\_features.json}, which lists both the feature names as they appear in the CICIDS2018 CSV headers and the order in which they are presented to the Random Forest and to the LLM agents. This ordering is significant because the structured natural language prompt used by each agent presents features in a fixed order, and any deviation would produce different tokenisation and potentially different reasoning paths.

Cost tracking is implemented at the per-agent level. Each agent records the number of input tokens, output tokens, and the USD cost of its API call using OpenAI's published pricing for the relevant model snapshot. These per-agent costs are aggregated to the per-flow level and then to the per-batch level in the result files. This granular cost recording enables post-hoc analysis of which agents are most expensive --- the Temporal Agent, which receives the longest context due to co-IP flow concatenation, is typically the most costly --- and supports the economic analysis presented in Chapter~6.

The pipeline control and logging infrastructure additionally supports reproducibility at the experimental level. The \texttt{run\_log.txt} file in the results directory records a timestamped entry for every batch start, checkpoint, and completion event, along with the software versions in use and the API model snapshot identifier at the time of execution. Should an experiment be interrupted and restarted, the log provides an unambiguous record of which flows had been processed prior to interruption, enabling the partial results to be understood in context rather than treated as a complete evaluation. All result files, batch files, and log files are committed to the project's version control repository upon completion of each experiment, establishing an immutable record of the experimental state at each stage of the research.
